\chapter{2002年北京卷（秋理）}
\section{单选题}
\begin{problem}
满足条件 \(M \cup \{1\} = \{1,2,3\}\) 的集合 \(M\) 的个数是（\quad）

\choices
    {\(1\)}
    {\(2\)}
    {\(3\)}
    {\(4\)}
\end{problem}

\begin{problem}
在平面直角坐标系中，已知两点 \(A(\cos 80^{\circ},\sin 80^{\circ})\)，\(B(\cos 20^{\circ},\sin 20^{\circ})\)，则 \(|AB|\) 的值是（\quad）

\choices
    {\(\frac{1}{2}\)}
    {\(\frac{\sqrt{2}}{2}\)}
    {\(\frac{\sqrt{3}}{2}\)}
    {\(1\)}
\end{problem}

\begin{problem}
下列四个函数中，以 \(\pi\) 为最小正周期，且在区间 \(\left(\frac{\pi}{2},\pi\right)\) 上为减函数的是（\quad）

\choices
    {\(y=\cos^{2}x\)}
    {\(y=2|\sin x|\)}
    {\(y=\left(\frac{1}{3}\right)^{\cos x}\)}
    {\(y=-\cot x\)}
\end{problem}

\begin{problem}
\(64\) 个直径都为 \(\frac{a}{4}\) 的球，记它们的体积之和为 \(V_{\text{甲}}\)，表面积之和为 \(S_{\text{甲}}\)；一个直径为 \(a\) 的球，记其体积为 \(V_{\text{乙}}\)，表面积为 \(S_{\text{乙}}\)，则（\quad）

\choices
    {\(V_{\text{甲}}>V_{\text{乙}},\ S_{\text{甲}}>S_{\text{乙}}\)}
    {\(V_{\text{甲}}<V_{\text{乙}}\)，\(\ S_{\text{甲}}<S_{\text{乙}}\)}
    {\(V_{\text{甲}}=V_{\text{乙}}\)，\(\ S_{\text{甲}}>S_{\text{乙}}\)}
    {\(V_{\text{甲}}=V_{\text{乙}}\)，\(\ S_{\text{甲}}=S_{\text{乙}}\)}
\end{problem}

\begin{problem}
已知某曲线的参数方程是 \(\begin{cases} x=\sec\varphi,\\ y=\tan\varphi,\end{cases}\)（\(\varphi\) 为参数），若以原点为极点，\(x\) 轴的正半轴为极轴，长度单位不变，建立极坐标系，则该曲线的极坐标方程是（\quad）

\choices
    {\(\rho=1\)}
    {\(\rho\cos 2\theta=1\)}
    {\(\rho^{2}\sin 2\theta=1\)}
    {\(\rho^{2}\cos 2\theta=1\)}
\end{problem}

\begin{problem}
给定四条曲线：

\begin{circlelist}
\item \(x^{2}+y^{2}=\frac{5}{2}\)
\item \(\frac{x^{2}}{9}+\frac{y^{2}}{4}=1\)
\item \(x^{2}+\frac{y^{2}}{4}=1\)
\item \(\frac{x^{2}}{4}+y^{2}=1\)
\end{circlelist}

其中与直线 \(x+y-\sqrt{5}=0\) 仅有一个交点的曲线是（\quad）

\choices
    {\(\circled{1}\circled{2}\circled{3}\)}
    {\(\circled{2}\circled{3}\circled{4}\)}
    {\(\circled{1}\circled{2}\circled{4}\)}
    {\(\circled{1}\circled{3}\circled{4}\)}
\end{problem}

\begin{problem}
已知 \(z_{1}\)，\(z_{2}\in\mathbb{C}\)，且 \(|z_{1}|=1\).若 \(z_{1}+z_{2}=2\)，则 \(|z_{1}-z_{2}|\) 的最大值是（\quad）

\choices
    {\(6\)}
    {\(5\)}
    {\(4\)}
    {\(3\)}
\end{problem}

\begin{problem}
若 \(\frac{\cot\theta-1}{2\cot\theta+1}=1\)，则 \(\frac{\cos 2\theta}{1+\sin 2\theta}\) 的值为（\quad）

\choices
    {\(3\)}
    {\(-3\)}
    {\(-2\)}
    {\(-\frac{1}{2}\)}
\end{problem}

\begin{problem}
\(12\) 名学生分别到三个不同的路口进行车流量的调查，若每个路口 \(4\) 人，则不同的分配方案共有（\quad）

\choices
    {\(\mathrm C_{12}^{4}\mathrm C_{8}^{4}\mathrm C_{4}^{4}\) 种}
    {\(3\mathrm C_{12}^{4}\mathrm C_{8}^{4}\mathrm C_{4}^{4}\) 种}
    {\(\mathrm C_{12}^{4}\mathrm C_{8}^{4}\mathrm A_{3}^{3}\) 种}
    {\(\frac{\mathrm C_{12}^{4}\mathrm C_{8}^{4}\mathrm C_{4}^{4}}{\mathrm A_{3}^{3}}\) 种}
\end{problem}

\begin{problem}
设命题甲：“直四棱柱 \(ABCD-A_{1}B_{1}C_{1}D_{1}\) 中，平面 \(ACB_{1}\) 与对角面 \(BB_{1}D_{1}D\) 垂直”；命题乙：“直四棱柱 \(ABCD-A_{1}B_{1}C_{1}D_{1}\) 是正方体”，那么甲是乙的（\quad）

\choices
    {充分必要条件}
    {充分非必要条件}
    {必要非充分条件}
    {即非充分又非必要条件}
\end{problem}

\begin{problem}
已知 \(f(x)\) 是定义在 \((-3,3)\) 上的奇函数，当 \(0<x<3\) 时，\(f(x)\) 的图象如图所示，那么不等式 \(f(x)\cos x<0\) 的解集是（\quad）

\begin{center}
\bitmapinclude[max width=0.42\linewidth]{img/2002/beijing_science_autumn/q11_fig1.png}
\end{center}

\choices
    {\((-3,-\frac{\pi}{2})\cup(0,1)\cup(\frac{\pi}{2},3)\)}
    {\((-\frac{\pi}{2},-1)\cup(0,1)\cup(\frac{\pi}{2},3)\)}
    {\((-3,-1)\cup(0,1)\cup(1,3)\)}
    {\((-3,-\frac{\pi}{2})\cup(0,1)\cup(1,3)\)}
\end{problem}

\begin{problem}
如图所示，\(f_i(x)\)（\(i=1\)，\(2\)，\(3\)，\(4\)）是定义在 \([0,1]\) 上的四个函数，其中满足性质：“对 \([0,1]\) 中任意的 \(x_1\) 和 \(x_2\)，任意 \(\lambda\in[0,1]\)，\(f[\lambda x_1+(1-\lambda)x_2]\le \lambda f(x_1)+(1-\lambda)f(x_2)\) 恒成立”的只有（\quad）

\begin{center}
\bitmapinclude[max width=0.95\linewidth]{img/2002/beijing_science_autumn/q12_fig1.png}
\end{center}

\choices
    {\(f_1(x),f_3(x)\)}
    {\(f_2(x)\)}
    {\(f_2(x)\)，\(f_3(x)\)}
    {\(f_4(x)\)}
\end{problem}

\section{填空题}
\begin{problem}
\(\arcsin\left(-\frac{2}{5}\right)\)，\(\arccos\left(-\frac{3}{4}\right)\)，\(\arctan\left(-\frac{5}{4}\right)\) 从小到大的顺序是\fillinblank{}.
\end{problem}

\begin{problem}
等差数列 \(\{a_n\}\) 中，\(a_1=2\)，公差不为零，且 \(a_1\)，\(a_3\)，\(a_{11}\) 恰好是某等比数列的前三项，那么该等比数列公比的值等于\fillinblank{}.
\end{problem}

\begin{problem}
关于直角 \(AOB\) 在平面 \(\alpha\) 内的射影有如下判断：

\begin{circlelist}
\item 可能是 \(0^{\circ}\) 的角
\item 可能是锐角
\item 可能是直角
\item 可能是直角
\item 可能是 \(180^{\circ}\) 的角
\end{circlelist}

其中正确的序号是\fillinblank{}.（注：把你认为正确判断的序号都填上）
\end{problem}

\begin{problem}
已知 \(P\) 是直线 \(3x+4y+8=0\) 上的动点，\(PA\)，\(PB\) 是圆 \(x^{2}+y^{2}-2x-2y+8=0\) 的两条切线，\(A\)，\(B\) 是切点，\(C\) 是圆心，那么四边形 \(PACB\) 面积的最小值为\fillinblank{}.
\end{problem}

\section{解答题}
\begin{problem}
解不等式 \(\left|\sqrt{2x-1}-x\right|<2\).
\end{problem}

\begin{problem}
如图，在多面体 \(ABCD-A_{1}B_{1}C_{1}D_{1}\) 中，上，下底面平行且均为矩形，相对的侧面与同一底面所成的二面角大小相等，侧棱延长后相交于 \(E\)，\(F\) 两点，上，下底面矩形的长，宽分别为 \(c\)，\(d\) 与 \(a\)，\(b\)，且 \(a>c\)，\(b>d\)，两底面间的距离为 \(h\).

\begin{enumerate}
\item 求侧面 \(ABB_{1}A_{1}\) 与底面 \(ABCD\) 所成二面角的大小；
\item 证明：\(EF \parallel ABCD\)；
\item 在估算该多面体的体积时，经常运用近似公式 \(V_{\text{估}}=S_{\text{中截面}}\cdot h\) 来计算，已知它的体积公式是 \(V=\frac{h}{6}(S_{\text{上底面}}+4S_{\text{中截面}}+S_{\text{下底面}})\)，试判断 \(V_{\text{估}}\) 与 \(V\) 的大小关系，并加以证明.
\end{enumerate}

\medskip
注：与两个底面平行，且到两个底面距离相等的截面称为该多面体的中截面.

\begin{flushright}
\bitmapinclude[max width=0.52\linewidth]{img/2002/beijing_science_autumn/q18_fig1.png}
\end{flushright}
\end{problem}

\begin{problem}
数列 \(\{x_n\}\) 由下列条件确定：\(x_1=a>0\)，\(x_{n+1}=\frac12\left(x_n+\frac{a}{x_n}\right)\)，\(n\in\N\).

\begin{enumerate}
\item 证明：对 \(n\ge 2\)，总有 \(x_n\ge \sqrt{a}\)；
\item 证明：对 \(n\ge 2\)，总有 \(x_n\ge x_{n+1}\)；
\item 若数列 \(\{x_n\}\) 的极限存在，且大于零，求 \(\lim\limits_{n\to\infty}x_n\) 的值.
\end{enumerate}
\end{problem}

\begin{problem}
在研究并行计算的基本算法时，有以下简单模型问题：用计算机求 \(n\) 个不同的数 \(v_1\)，\(v_2\)，\(\cdots\)，\(v_n\) 的和 \(\sum_{i=1}^{n}v_i=v_1+v_2+\cdots+v_n\)，计算开始前，\(n\) 个数存贮在 \(n\) 台由网络连接的计算机中，每台机器存一个数. 计算开始后，在一个单位时间内，每台机器至多到一台其他机器中读数据，并与自己原有数据相加得到新的数据，各台机器可同时完成上述工作. 为了用尽可能少的单位时间，即可完成计算，方法可用下表表示：

\begin{center}
\begingroup
\setlength{\tabcolsep}{4pt}
\renewcommand{\arraystretch}{1.15}
\begin{tabular}{|c|c|cc|cc|cc|}
\hline
\multirow{2}{*}{机器号} & \multirow{2}{*}{初始时} & \multicolumn{2}{c|}{第一单位时间} & \multicolumn{2}{c|}{第二单位时间} & \multicolumn{2}{c|}{第三单位时间} \\
\cline{3-8}
 & & 被读机号 & 结果 & 被读机号 & 结果 & 被读机号 & 结果 \\
\hline
1 & \(v_1\) & 2 & \(v_1+v_2\) &  &  &  &  \\
\hline
2 & \(v_2\) & 1 & \(v_2+v_1\) &  &  &  &  \\
\hline
\end{tabular}
\endgroup
\end{center}

\begin{enumerate}
\item 当 \(n=4\) 时，至少需要多少个单位时间可完成计算? 把你设计的方法填入下表：

\begin{center}
\begingroup
\setlength{\tabcolsep}{4pt}
\renewcommand{\arraystretch}{1.15}
\begin{tabular}{|c|c|cc|cc|cc|}
\hline
\multirow{2}{*}{机器号} & \multirow{2}{*}{初始时} & \multicolumn{2}{c|}{第一单位时间} & \multicolumn{2}{c|}{第二单位时间} & \multicolumn{2}{c|}{第三单位时间} \\
\cline{3-8}
 & & 被读机号 & 结果 & 被读机号 & 结果 & 被读机号 & 结果 \\
\hline
1 & \(v_1\) & & & & & & \\
\hline
2 & \(v_2\) & & & & & & \\
\hline
3 & \(v_3\) & & & & & & \\
\hline
4 & \(v_4\) & & & & & & \\
\hline
\end{tabular}
\endgroup
\end{center}

\item 当 \(n=128\) 时，要使所有机器都得到 \(\sum_{i=1}^{n}v_i\)，至少需要多少个单位时间可完成计算?（结论不要求证明）
\end{enumerate}
\end{problem}

\begin{problem}
已知 \(O(0,0)\)，\(B(1,0)\)，\(C(b,c)\) 是 \(\triangle OBC\) 的三个顶点.

\begin{enumerate}
\item 写出 \(\triangle OBC\) 的重心 \(G\)，外心 \(F\)，垂心 \(H\) 的坐标，并证明 \(G\)，\(F\)，\(H\) 三点共线；
\item 当直线 \(FH\) 与 \(OB\) 平行时，求顶点 \(C\) 的轨迹.
\end{enumerate}
\end{problem}

\begin{problem}
已知 \(f(x)\) 是定义在 \(\R\) 上的不恒为零的函数，且对于任意的 \(a\)，\(b\in\R\) 都满足：\(f(a\cdot b)=af(b)+bf(a)\).

\begin{enumerate}
\item 求 \(f(0)\)，\(f(1)\) 的值；
\item 判断 \(f(x)\) 的奇偶性，并证明你的结论；
\item 若 \(f(2)=2\)，\(u_n=\frac{f(2^{-n})}{n}\)（\(n\in\N\)），求数列 \(\{u_n\}\) 的前 \(n\) 项的和 \(S_n\).
\end{enumerate}
\end{problem}
